\Askhsh{25745_SOLUTION}{1}
\begin{ylh}{1}
    \item [Ύλη από εικόνα] Λύση
\end{ylh}
\begin{alist}
    \item[α)]
    \begin{rlist}
        \item [i.] Αν υπήρχε $x_o \in (0,2)$ με $f(x_o)=0$ τότε από τη δοσμένη σχέση θα είχαμε $(f'(x_o))^2 < 0$ το οποίο είναι άτοπο. Συνεπώς $f(x) \neq 0$ για κάθε $x \in (0,2)$.
        \item [ii.] Αφού $f$ συνεχής στο $(0,2)$ και $f(x) \neq 0$ για κάθε $x \in (0,2)$ θα διατηρεί πρόσημο σε αυτό και επειδή $f(1)=1>0$ θα είναι $f(x)>0$ για κάθε $x \in (0,2)$.
    \end{rlist}
    \item[β)] Από τη δοσμένη σχέση και αφού δείξαμε ότι $f(x)>0$ για κάθε $x \in (0,2)$ προκύπτει ότι $f''(x) < -\frac{(f'(x))^2}{f(x)} < 0$ για κάθε $x \in (0,2)$, οπότε αφού η $f$ είναι συνεχής στο $[0,2]$ είναι κοίλη στο $[0,2]$ και δεν έχει σημεία καμπής.
    \item[γ)] Αφού $f$ κοίλη στο $[0,2]$ η $f'$ θα είναι γνησίως φθίνουσα στο $(0,2)$ οπότε για κάθε $0 < x < 1 \xrightarrow{f'\downarrow} f'(x) > f'(1) \Rightarrow f'(x) > 0$ και αφού $f$ συνεχής στο $[0,1]$ θα είναι γνησίως αύξουσα στο $[0,1]$, ενώ για κάθε $1 < x < 2 \xrightarrow{f'\downarrow} f'(x) < f'(1) \Rightarrow f'(x) < 0$ και αφού $f$ συνεχής στο $[1,2]$ θα είναι γνησίως φθίνουσα στο $[1,2]$.
    Συνεπώς η $f$ παρουσιάζει ολικό μέγιστο στο $1$ και ολικό ελάχιστο στο $0$ και στο $2$, αφού ισχύει $f(0)=f(2)$.
\end{alist}

\noindent
\textbf{Σημείωση:} Μία συνάρτηση που ικανοποιεί τις υποθέσεις του θέματος είναι η $f(x) = \sqrt{-x^2+2x}$ για κάθε $x \in [0,2]$.